% chapters/02/latihan.tex
\bagian{Latihan dan Refleksi}

\begin{enumerate}
  \item Sebutkan empat unsur pembentuk permainan, lalu jelaskan masing-masing
        dengan contoh dari satu permainan yang Anda kenal.

  \item Jelaskan mengapa program permainan digolongkan sebagai sistem waktu
        nyata, dan berikan satu contoh akibat yang timbul apabila batas
        waktunya dilanggar.

  \item Sebuah objek bergerak dengan kecepatan 150 piksel per detik. Hitunglah
        perpindahannya pada satu bingkai apabila laju bingkainya 50 FPS, lalu
        jelaskan apa yang terjadi pada hasil perhitungan apabila selang waktu
        tidak dipakai sama sekali.

  \item Jelaskan perbedaan langkah tetap dan langkah berubah. Sebutkan satu
        jenis permainan yang lebih cocok memakai masing-masing cara.

  \item Mengapa arah sumbu tegak pada koordinat layar berbeda dengan koordinat
        kartesius, dan apa akibatnya ketika Anda menuliskan pengaruh gravitasi
        pada kode program?

  \item Jelaskan mengapa galat kaidah Euler menyusut ketika panjang langkah
        diperkecil. Sebutkan pula dua cara lain untuk menjaga kestabilan
        simulasi selain memperkecil panjang langkah.

  \item Jelaskan perbedaan lapisan keputusan dan lapisan navigasi pada
        pengendalian NPC beserta isi jawaban masing-masing. Berikan satu contoh
        keadaan yang hanya dapat ditangani oleh lapisan navigasi.

  \item Sebutkan tiga pekerjaan yang diambil alih oleh mesin permainan, lalu
        jelaskan satu keuntungan dan satu kerugian dari pengambilalihan itu.

  \item Perhatikan Tabel~\ref{tab:padanan-konsep}. Pilih satu baris, lalu
        jelaskan mengapa konsep tersebut tetap sama meskipun cara
        penerapannya berbeda pada kedua perkakas.
\end{enumerate}

\begin{refleksibox}
\begin{enumerate}
  \item Setelah membaca bab ini, bagian mana yang paling mengubah cara Anda
        memandang permainan yang biasa Anda mainkan?
  \item Adakah istilah pada bab ini yang masih terasa asing? Catatlah istilah
        tersebut, lalu bandingkan catatan Anda dengan Glosarium pada
        Lampiran~C sebelum pertemuan berikutnya.
  \item Pada Telaah~3, keadaan NPC mana yang paling sulit Anda kenali?
        Mengapa demikian?
\end{enumerate}
\end{refleksibox}
